% =============================================================================
\chapter{Cài đặt và kiểm thử}\label{chap:caidat}
% =============================================================================

\section{Môi trường và cách biên dịch}

Chương trình viết bằng C++17, biên dịch bằng \code{g++} 13.3 với các cờ \code{-std=c++17 -O2 -Wall -Wextra -pedantic}
(không có cảnh báo nào). Mã nguồn không phụ thuộc hệ điều hành và biên dịch được bằng MinGW trên Windows.
Tệp \code{Makefile} cung cấp các lệnh:

\begin{lstlisting}[language=bash,caption={Các lệnh biên dịch và chạy},label=lst:make]
make              # bien dich chuong trinh chinh ./quanlykho
make run          # chay menu voi du lieu trong data/
make test         # bien dich va chay toan bo test
make bench        # thuc nghiem B-Tree, ghi ketqua/ketqua_btree.csv
./quanlykho data 3        # tham so: thu muc du lieu, bac B-Tree (mac dinh: data 5)
\end{lstlisting}

\section{Cấu trúc mã nguồn}

\begin{table}[htbp]
\centering\small
\caption{Các tệp mã nguồn (số dòng, kể cả chú thích)}
\label{tab:tep_nguon}
\begin{tabularx}{\linewidth}{@{}l r Y l@{}}
\toprule
\textbf{Mô-đun} & \textbf{Dòng} & \textbf{Nội dung chính} & \textbf{Phụ trách} \\
\midrule
\code{Types.h} & 59 & struct dùng chung, \code{enum KetQua} & Cả nhóm \\
\code{BTree.h/.cpp} & 363 & tìm, chèn/tách, xóa/mượn/gộp, duyệt, in cây & \TVa \\
\code{Benchmark.h/.cpp} & 160 & thực nghiệm $m = 3, 5, 10$, ghi CSV & \TVa, \TVd \\
\code{LotQueue.h/.cpp} & 90 & hàng đợi lô (DSLK) & \TVb \\
\code{Phieu.h/.cpp} & 111 & danh sách phiếu (DSLK), lọc theo ngày & \TVb \\
\code{Inventory.h/.cpp} & 239 & thêm/sửa/xóa mặt hàng, nhập kho, xuất kho FIFO & \TVb \\
\code{Sort.h/.cpp} & 135 & mảng động, Merge Sort, các tiêu chí sắp xếp & \TVc \\
\code{Report.h/.cpp} & 281 & tìm kiếm, cảnh báo, báo cáo, in bảng & \TVc \\
\code{Date.h/.cpp} & 118 & kiểm tra, quy đổi, cộng/trừ/so sánh ngày & \TVd \\
\code{Validate.h/.cpp} & 154 & kiểm tra chuỗi, nhập có kiểm tra & \TVd \\
\code{FileIO.h/.cpp} & 329 & quy đổi mã, đọc/ghi ba tệp & \TVd \\
\code{DataGen.h/.cpp} & 139 & sinh dữ liệu mẫu và dữ liệu lớn & \TVd \\
\code{main.cpp} & 360 & menu 18 chức năng & \TVd \\
\midrule
\textbf{Tổng mã nguồn} & \textbf{2538} & & \\
Mã kiểm thử (\code{tests/}) & 1744 & 7 bộ kiểm thử & \\
\bottomrule
\end{tabularx}
\end{table}

\section{Giao diện chương trình}\label{sec:manhinh}

Chương trình chạy trên cửa sổ dòng lệnh. Khi khởi động, chương trình đọc ba tệp dữ liệu và hiện menu chính
(Hình~\ref{fig:menu}); dòng thứ hai cho biết số mặt hàng, lô, phiếu, bậc và chiều cao B-Tree, thư mục dữ liệu và
trạng thái chưa lưu. Các hình trong mục này là kết quả chạy thật của chương trình trên bộ dữ liệu mẫu
\code{data/} ngày 26/11/2026; phần người dùng gõ được terminal hiển thị lại sau dấu nhắc. Để gọn, menu chỉ được
in ở hình đầu tiên.

\begin{figure}[htbp]
\begin{manhinh}[./quanlykho data]
Dang doc du lieu tu data/ ...
Da nap 35 mat hang, 65 lo, 99 phieu.

==================== QUAN LY KHO HANG SIEU THI ====================
 Kho: 35 mat hang, 65 lo, 99 phieu | B-Tree m = 5, cao 3 | data/
 MAT HANG   1. Them   2. Sua   3. Xoa   4. Xem chi tiet theo ma vach
            5. Tim theo ten   6. Liet ke theo nhom
 NHAP/XUAT  7. Nhap kho   8. Xuat kho (FIFO)   9. Danh sach phieu
 BAO CAO    10. Ton duoi nguong   11. Sap het han (7 ngay)
            12. Gia tri ton theo nhom   13. Top 10 ban chay trong thang
 B-TREE     14. Xem cay   15. Bat/tat in cay khi tach/gop   16. Thuc nghiem m = 3, 5, 10
 DU LIEU    17. Sinh du lieu mau   18. Luu   0. Luu va thoat
===================================================================
Chon: 99
  -> Loi: phai nam trong [0, 18]. Nhap lai.
Chon: abc
  -> Loi: phai la so nguyen. Nhap lai.
\end{manhinh}
\caption{Khởi động và menu chính; nhập lựa chọn sai thì được hỏi lại}
\label{fig:menu}
\end{figure}

\paragraph{Xem chi tiết và xuất kho FIFO.} Mục 4 tìm mặt hàng qua B-Tree và cho biết số nút đã truy xuất.
Mục 8 xuất 30 đơn vị: lô 4 (25 đơn vị) bị lấy hết và gỡ khỏi hàng đợi, 5 đơn vị còn lại lấy từ lô 5
(Hình~\ref{fig:xuatkho}).

\begin{figure}[htbp]
\begin{manhinh}[Muc 4 va muc 8]
Chon: 4
Ma vach: 8930000000002
---------------- THONG TIN MAT HANG ----------------
Ma vach         : 8930000000002
Ten mat hang    : Mi goi Loai 2
Nhom hang       : Thuc pham kho
Tong ton kho    : 64
So lo hien co   : 2
Chi tiet cac lo hang (FIFO - Cu nhat den moi nhat):
  Ma Lo   Ngay Nhap     Han Su Dung   So Luong  Gia Nhap
  4       09/11/2026    15/04/2027    25        82500.00
  5       16/11/2026    18/04/2027    39        82500.00
----------------------------------------------------
(B-Tree: tim qua 3 nut, cay cao 3)

Chon: 8
Ma vach: 8930000000002
   ... (in lai thong tin mat hang nhu tren) ...
So luong xuat (1 - 64): 30
Ngay xuat dd/mm/yyyy (Enter = 26/11/2026):
Da tao phieu xuat PX00100 (lay tu lo cu nhat truoc).
   ...
Tong ton kho    : 34
So lo hien co   : 1
Chi tiet cac lo hang (FIFO - Cu nhat den moi nhat):
  Ma Lo   Ngay Nhap     Han Su Dung   So Luong  Gia Nhap
  5       16/11/2026    18/04/2027    34        82500.00
----------------------------------------------------
\end{manhinh}
\caption{Xem chi tiết mặt hàng (kèm số nút B-Tree truy xuất) và xuất kho theo FIFO}
\label{fig:xuatkho}
\end{figure}

\paragraph{Nhập kho với kiểm tra dữ liệu.} Mọi giá trị nhập đều được kiểm tra kiểu và miền; nhập sai thì báo
lý do cụ thể và hỏi lại, bấm Enter ở ô ngày để lấy ngày hôm nay (Hình~\ref{fig:nhapkho}).

\begin{figure}[htbp]
\begin{manhinh}[Muc 7 - Nhap kho]
Chon: 7
Ma vach: 8930000000002
Mi goi Loai 2 - ton hien tai: 34
So luong nhap: 0
  -> Loi: phai nam trong [1, 1000000]. Nhap lai.
So luong nhap: 20
Gia nhap: 85000
Ngay nhap dd/mm/yyyy (Enter = 26/11/2026):
Han su dung dd/mm/yyyy: 31/02/2027
  -> Loi: ngay khong hop le (dd/mm/yyyy, nam 1900-2100). Nhap lai.
Han su dung dd/mm/yyyy: 01/10/2026
  -> Loi: HSD khong duoc truoc ngay nhap 26/11/2026. Nhap lai.
Han su dung dd/mm/yyyy: 15/06/2027
Da tao phieu nhap PN00101, lo LO00072. Ton moi: 54
\end{manhinh}
\caption{Nhập kho: số lượng ngoài miền, ngày không tồn tại, HSD trước ngày nhập đều bị từ chối}
\label{fig:nhapkho}
\end{figure}

\paragraph{Cảnh báo và báo cáo.} Hình~\ref{fig:canhbao} là hai cảnh báo: tồn dưới ngưỡng (sắp theo tồn tăng
dần) và lô hết hạn trong 7 ngày (sắp theo HSD, gồm cả lô đã quá hạn mà vẫn còn hàng). Hình~\ref{fig:baocao}
là báo cáo giá trị tồn theo nhóm và top 10 bán chạy tháng 11/2026. Mì gói Loại 2 đứng đầu với
$51 + 30 = 81$ đơn vị: 51 từ phiếu \code{PX00006} trong dữ liệu mẫu và 30 từ lần xuất ở Hình~\ref{fig:xuatkho}.

\begin{figure}[htbp]
\begin{manhinh}[Muc 10 va muc 11]
Chon: 10
Mat hang co ton kho DUOI nguong canh bao:
STT  Ma vach        Ten                       Nhom                 Don gia     Ton  Nguong
------------------------------------------------------------------------------------------
1    8930000000018  Gao Loai 18               Thuc pham kho          55000      14      22
2    8930000000021  Ca phe Loai 21            Do uong               105000      16      20
Tong: 2 mat hang

Chon: 11
Lo con hang da het han hoac het han trong 7 ngay (tinh tu 26/11/2026):
STT  Ma vach        Ten                       Ma lo   HSD         Con lai             SL
----------------------------------------------------------------------------------------
1    8930000000030  Sua tuoi Loai 30          60      23/11/2026  DA HET HAN          28
2    8930000000025  Dau an Loai 25            49      24/11/2026  DA HET HAN          28
3    8930000000020  Tra xanh Loai 20          39      25/11/2026  DA HET HAN          21
4    8930000000015  Sua tuoi Loai 15          30      26/11/2026  0 ngay              31
5    8930000000010  Dau an Loai 10            19      27/11/2026  1 ngay               8
\end{manhinh}
\caption{Cảnh báo tồn dưới ngưỡng và cảnh báo hạn sử dụng}
\label{fig:canhbao}
\end{figure}

\begin{figure}[htbp]
\begin{manhinh}[Muc 12 va muc 13]
Chon: 12
Gia tri ton kho theo nhom (tinh theo gia nhap tung lo):
Nhom                   So MH    Tong ton       Gia tri ton
----------------------------------------------------------
Thuc pham tuoi             5         265          34916250
Hoa my pham                4         263          32298750
Do uong                    8         351          29936250
Do gia dung                4         242          23977500
Thuc pham kho              6         278          19647500
Banh keo                   4         247          15311250
Gia vi                     4         134           4837500
----------------------------------------------------------
Tong gia tri ton (theo gia nhap)                 160925000

Chon: 13
Thang (Enter = 11):
Nam (Enter = 2026):
Top 10 ban chay thang 11/2026 (theo phieu xuat):
Hang  Ma vach        Ten                       Nhom                SL ban
-------------------------------------------------------------------------
1     8930000000002  Mi goi Loai 2             Thuc pham kho           81
2     8930000000005  Tra xanh Loai 5           Do uong                 54
3     8930000000014  Nuoc rua chen Loai 14     Do gia dung             49
4     8930000000032  Mi goi Loai 32            Thuc pham kho           49
5     8930000000035  Tra xanh Loai 35          Do uong                 43
6     8930000000023  Keo deo Loai 23           Banh keo                39
7     8930000000004  Nuoc ngot Loai 4          Do uong                 33
8     8930000000017  Mi goi Loai 17            Thuc pham kho           31
9     8930000000010  Dau an Loai 10            Gia vi                  30
10    8930000000008  Keo deo Loai 8            Banh keo                24
\end{manhinh}
\caption{Báo cáo giá trị tồn theo nhóm và top 10 bán chạy trong tháng}
\label{fig:baocao}
\end{figure}

\paragraph{In cây khi tách/gộp (nâng cao).} Mục 15 bật chế độ in cây. Hình~\ref{fig:intachgop} chạy với kho rỗng
và B-Tree bậc 3: mặt hàng thứ ba làm lá tràn và bị tách; xóa một mặt hàng làm lá rỗng và hai nút gộp lại,
cây thấp đi một mức.

\begin{figure}[htbp]
\begin{manhinh}[./quanlykho nho 3]
Chon: 15
In cay sau moi lan tach/muon/gop nut: BAT
Chon: 1
Ma vach (8 hoac 13 chu so): 00000080
Ten mat hang: Nuoc ngot
   ... (da them 00000050 va 00000020 truoc do) ...
  Tach nut: day 00000050 len, [00000020] | [00000080]
Muc 0: [00000050]
Muc 1: [00000020] [00000080]
Thanh cong.

Chon: 3
Ma vach: 00000080
   ... (thong tin mat hang) ...
Xac nhan xoa mat hang nay? (c/k): c
  Gop nut voi khoa 00000050 -> [00000020 00000050]
Muc 0: [00000020 00000050]
Thanh cong.
\end{manhinh}
\caption{In cây sau mỗi lần tách và gộp nút (B-Tree bậc 3)}
\label{fig:intachgop}
\end{figure}

\paragraph{Kiểm tra dữ liệu khi đọc tệp.} Nếu tệp dữ liệu có dòng sai (sửa tay hoặc hỏng), chương trình bỏ qua
dòng đó, báo vị trí và lý do, và vẫn nạp phần dữ liệu đúng (Hình~\ref{fig:doctep}).

\begin{figure}[htbp]
\begin{manhinh}[./quanlykho loi]
Dang doc du lieu tu loi/ ...
  [mathang.txt dong 2] ma vach phai co 8 hoac 13 chu so -> bo qua: 893000000002|Ma 12 so|Do uong|5000|5
  [mathang.txt dong 3] trung ma vach -> bo qua: 8930000000001|Trung ma|Do uong|5000|5
  [mathang.txt dong 4] don gia sai -> bo qua: 8930000000003|Gia am|Do uong|-500|5
  [mathang.txt dong 5] don gia sai -> bo qua: 8930000000004|Keo deo|Banh keo|abc|5
  [lo.txt dong 2] ngay sai -> bo qua: 8930000000001|LO00002|30/02/2026|01/12/2026|5|20000
  [lo.txt dong 3] HSD truoc ngay nhap -> bo qua: 8930000000001|LO00003|10/11/2026|01/11/2026|5|20000
  [lo.txt dong 4] ma vach khong co trong mathang.txt -> bo qua: 8930000000009|LO00004|...
  [phieu.txt dong 2] ma phieu sai (PN/PX + so, khop loai) -> bo qua: PX00002|NHAP|05/11/2026|...
Mot so dong loi da bi bo qua (xem thong bao o tren).
Da nap 1 mat hang, 1 lo, 1 phieu.
\end{manhinh}
\caption{Đọc tệp có dòng lỗi: mã vạch sai, trùng mã, sai kiểu, sai miền, sai tiền tố mã phiếu}
\label{fig:doctep}
\end{figure}

\section{Chiến lược kiểm thử}

Mỗi mô-đun có một chương trình kiểm thử tự động trong thư mục \code{tests/}; lệnh \code{make test} chạy tất cả.
Các test gồm ba loại theo yêu cầu: \textbf{bình thường}, \textbf{biên} và \textbf{dữ liệu sai}. Ngoài các ca
cụ thể, nhóm dùng thêm các kỹ thuật:

\begin{itemize}
  \item \textbf{Kiểm tra bất biến B-Tree} sau \emph{mỗi} thao tác chèn/xóa: số khóa mỗi nút trong
        $[t-1, m-1]$, khóa tăng dần và nằm đúng khoảng của cha, mọi lá cùng độ sâu, tổng số khóa bằng \code{size}.
  \item \textbf{Kiểm thử ngẫu nhiên đối chiếu mô hình}: 20\,000 thao tác thêm/nhập/xuất/xóa ngẫu nhiên trên
        2\,000 mã vạch với B-Tree bậc 3; sau mỗi bước so danh sách lô và tồn kho với một mô hình FIFO đơn giản viết
        bằng mảng.
  \item \textbf{Đọc rồi ghi lại}: đọc bộ dữ liệu rồi ghi ra phải được ba tệp giống hệt từng byte; thao tác, lưu,
        đọc lại (với bậc B-Tree khác) phải ra đúng trạng thái cũ.
  \item \textbf{Chạy menu bằng kịch bản}: \code{tests/test\_menu.sh} đưa sẵn chuỗi lựa chọn (kể cả nhập sai) vào
        chương trình thật trên bản sao của \code{data/}, rồi kiểm tra 36 điều kiện: 29 trên màn hình, 7 trên tệp đã lưu và mã thoát.
  \item \textbf{Kiểm tra lỗi bộ nhớ và rò rỉ}: toàn bộ test được biên dịch lại với AddressSanitizer,
        LeakSanitizer và UndefinedBehaviorSanitizer (\code{-fsanitize=address,undefined}); công cụ báo lỗi ngay khi
        truy cập vùng nhớ sai, giải phóng hai lần hoặc còn vùng nhớ chưa giải phóng lúc thoát.
\end{itemize}

\section{Các ca kiểm thử}

Ký hiệu loại: \textbf{BT} -- bình thường, \textbf{Biên} -- giá trị biên, \textbf{Sai} -- dữ liệu sai.
Tất cả các ca dưới đây đều đạt (\ok).

\begingroup\small
\begin{longtable}{@{}c c >{\raggedright\arraybackslash}p{6.2cm} >{\raggedright\arraybackslash}p{5.6cm} c@{}}
\caption{Các ca kiểm thử}\label{tab:testcase}\\
\toprule
\textbf{STT} & \textbf{Loại} & \textbf{Đầu vào} & \textbf{Kết quả mong đợi} & \\
\midrule
\endfirsthead
\toprule
\textbf{STT} & \textbf{Loại} & \textbf{Đầu vào} & \textbf{Kết quả mong đợi} & \\
\midrule
\endhead
\multicolumn{5}{@{}l}{\textit{B-Tree (\code{test\_btree})}} \\
1 & BT & Chèn 2\,000 khóa theo thứ tự ngẫu nhiên, $m \in \{3,4,5,10\}$ & Bất biến đúng sau mỗi lần chèn; tìm thấy mọi khóa & \ok \\
2 & BT & Duyệt cây sau khi chèn & Đủ 2\,000 khóa, tăng dần & \ok \\
3 & Biên & Chèn 500 khóa tăng dần rồi xóa giảm dần & Bất biến đúng sau mỗi bước; cây rỗng & \ok \\
4 & Biên & Tìm, xóa, chèn \code{nullptr} trên cây rỗng & Trả \code{nullptr}/\textbf{sai}; chiều cao 0 & \ok \\
5 & Sai & Chèn mã vạch đã có & Trả \textbf{sai}, cây không đổi & \ok \\
6 & Sai & Xóa khóa đã xóa & Trả \code{nullptr} & \ok \\
7 & BT & $10^5$ khóa mỗi bậc: chèn, tìm, xóa một nửa, xóa hết & Khóa đã xóa không còn, khóa còn lại tìm được; cây rỗng & \ok \\
8 & BT & Demo $m=3$: chèn 10 khóa, xóa 5 khóa & Đúng các bước ở Hình~\ref{fig:chen}, \ref{fig:xoa} & \ok \\
\midrule
\multicolumn{5}{@{}l}{\textit{Hàng đợi lô, phiếu, nhập/xuất kho (\code{test\_kho})}} \\
9 & BT & Nhập 3 lô 10, 20, 30 & Tồn 60; 3 lô đúng thứ tự; mã phiếu 1, 2, 3 & \ok \\
10 & BT & Xuất 15 khi các lô là $[9, 20, 30, 1]$ & Lô 1 bị gỡ, lô 2 còn 14, tồn 45 & \ok \\
11 & Biên & Xuất đúng bằng số lượng lô đầu (14) & Lô bị gỡ, \code{front} sang lô kế & \ok \\
12 & Biên & Xuất hết sạch; xuất tiếp 1 & Hàng đợi rỗng, \code{front = rear = nullptr}; \code{KQ\_KHONG\_DU\_HANG} & \ok \\
13 & Sai & Xuất 61 khi tồn 60 & \code{KQ\_KHONG\_DU\_HANG}; tồn, lô, phiếu không đổi & \ok \\
14 & Sai & Nhập số lượng 0/âm, giá âm, HSD trước ngày nhập, ngày 31/11, 30/02, năm 1899, 2101 & \code{KQ\_DU\_LIEU\_SAI}; không tạo lô, phiếu & \ok \\
15 & Sai & Thêm mặt hàng: 7 dạng mã vạch sai, tên/nhóm rỗng, giá âm/NaN/$\infty$, ngưỡng âm & \code{KQ\_DU\_LIEU\_SAI} & \ok \\
16 & Sai & Thêm mã vạch trùng & \code{KQ\_TRUNG\_MA} & \ok \\
17 & Biên & Xóa mặt hàng còn tồn; xóa kèm tồn kho & \code{KQ\_CON\_TON\_KHO}; sau đó \code{KQ\_OK} & \ok \\
18 & Biên & Nhập thêm \code{INT\_MAX} đơn vị & \code{KQ\_DU\_LIEU\_SAI} (chống tràn số) & \ok \\
19 & BT & \code{dspThemCoSan} với mã 50 rồi mã 7 & Mã kế tiếp 51, không giảm & \ok \\
20 & BT & 20\,000 thao tác ngẫu nhiên, đối chiếu mô hình FIFO & Lô và tồn kho khớp sau mỗi bước & \ok \\
\midrule
\multicolumn{5}{@{}l}{\textit{Tìm kiếm, sắp xếp, báo cáo (\code{test\_report})}} \\
21 & BT & Merge Sort mảng $n = 0, 1, 2, 3, 7, 16, 100, 10^4$ có khóa trùng & Đúng thứ tự, ổn định & \ok \\
22 & BT & Tìm \code{"sua"}, \code{"SUA T"} & Không phân biệt hoa thường, sắp theo tên & \ok \\
23 & Biên & Tìm \code{"xyz"}; từ khóa rỗng & Rỗng; tất cả mặt hàng & \ok \\
24 & Biên & Nhóm \code{"do uong"} khi có \code{"Do uong"} và \code{"DO UONG"} & Gộp thành một nhóm & \ok \\
25 & Biên & Tồn đúng bằng ngưỡng; ngưỡng 0 & Không cảnh báo & \ok \\
26 & Biên & Lô có HSD còn $-1, 0, 3, 7, 8$ ngày & Lấy 4 lô $-1..7$, sắp theo HSD; không kèm hết hạn: 3 lô & \ok \\
27 & BT & Giá trị tồn theo nhóm, tính tay & 290\,000; 146\,000; 0 & \ok \\
28 & Biên & Top bán chạy: phiếu tháng khác, phiếu nhập, gọi hai lần, $k = 0$, $k = 2$ & Chỉ tính phiếu xuất trong tháng; không cộng dồn & \ok \\
29 & Sai & Phiếu xuất của mặt hàng đã xóa & Bị bỏ qua & \ok \\
30 & BT & $10^5$ mặt hàng: tìm \code{"pham 99"}, đếm lô sắp hết hạn & 1\,111 và 11\,669 (đúng công thức) & \ok \\
\midrule
\multicolumn{5}{@{}l}{\textit{Ngày và kiểm tra nhập (\code{test\_date}, \code{test\_fileio})}} \\
31 & BT & Mọi ngày 01/01/1900 -- 31/12/2100 & 73\,414 ngày; đổi qua lại, định dạng, \code{congNgay} đúng & \ok \\
32 & Biên & Năm nhuận 2000, 2024; không nhuận 1900, 2026 & 29/02 hợp lệ đúng năm & \ok \\
33 & Sai & 17 chuỗi ngày sai: \code{29/02/2025}, \code{31/04/2026}, \code{01-01-2026}, \code{01/01/2026x}, \ldots & Từ chối, không ghi đè kết quả & \ok \\
34 & Sai & Số: \code{"2147483648"}, \code{"1.5"}, \code{"0x10"}, \code{"1,5"}, \code{"nan"}, \code{"1e999"} & Từ chối & \ok \\
35 & Sai & Nhập \code{abc}, \code{500}, \code{-1}, \code{42} với miền $[0, 100]$ & Hỏi lại 3 lần, nhận 42 & \ok \\
36 & Biên & Hết dữ liệu nhập giữa chừng & Lưu dữ liệu và thoát, không lặp vô hạn & \ok \\
\midrule
\multicolumn{5}{@{}l}{\textit{Đọc/ghi tệp (\code{test\_fileio}, \code{test\_dulieu}, \code{test\_menu.sh})}} \\
37 & BT & Đọc \code{data/} & 35 mặt hàng, 65 lô, 99 phiếu; tồn = tổng lô & \ok \\
38 & BT & Đọc rồi ghi lại; $10^4$ mặt hàng & Ba tệp giống hệt từng byte & \ok \\
39 & BT & Thêm, nhập, xuất, xóa; lưu; đọc lại với bậc khác & Trạng thái giống hệt; giữ số thập phân & \ok \\
40 & Sai & Tệp 17 dòng lỗi (sai số trường, mã, ngày, FIFO, mã phiếu trùng\ldots) & Bỏ qua đúng 17 dòng, báo số dòng; khi lưu, tệp gốc giữ nguyên trong \code{.bak} & \ok \\
41 & Biên & Tệp CRLF, dòng trống; thiếu \code{lo.txt}; thư mục không tồn tại & Đọc bình thường; vẫn nạp mặt hàng; kho rỗng & \ok \\
42 & BT & Tính nhất quán \code{data/}: lô khớp phiếu, thứ tự FIFO, mã lô/phiếu & 0 lỗi, 0 cảnh báo & \ok \\
43 & BT & Kịch bản menu 18 chức năng có nhập sai xen kẽ & 36 điều kiện trên màn hình, tệp và mã thoát đều đúng & \ok \\
\bottomrule
\end{longtable}
\endgroup

\section{Kết quả chạy kiểm thử và kiểm tra bộ nhớ}

Hình~\ref{fig:maketest} trích các dòng tổng kết khi chạy \code{make test}. Toàn bộ test cũng được chạy lại với
AddressSanitizer, LeakSanitizer và UndefinedBehaviorSanitizer (lệnh dưới đây): không có báo lỗi truy cập
bộ nhớ, hành vi không xác định hay rò rỉ bộ nhớ nào.
\begin{lstlisting}[language=bash,numbers=none]
make clean test CXXFLAGS="-std=c++17 -g -O1 -fsanitize=address,undefined"
\end{lstlisting} Khi thoát, chương trình gọi \code{khoGiaiPhong}: duyệt cây để giải
phóng từng mặt hàng cùng các lô, rồi giải phóng các nút cây và danh sách phiếu.

\begin{figure}[htbp]
\begin{manhinh}[make test (trich dong tong ket)]
TAT CA TEST DEU QUA
TAT CA TEST DATE DEU QUA
  con 1892 mat hang, 9015 phieu, cay cao 9
TAT CA TEST KHO DEU QUA
  tim 'pham 99': 1111, ton thap: 46667, sap het han: 11669, nhom: 5, top: 10 (131 ms)
TAT CA TEST REPORT DEU QUA
  doc 56 ms, ghi 38 ms (28000 phieu)
TAT CA TEST FILEIO DEU QUA
Ham quy doi maLo/maPhieu: dung
Ket qua: 0 LOI, 0 CANH BAO
TAT CA TEST MENU DEU QUA
\end{manhinh}
\caption{Kết quả \code{make test}}
\label{fig:maketest}
\end{figure}
